\begin{definition}{Ordnungsstatistik}{Ordnungsstatistik}
Seien $n\in \N$, $k \in \haken n$ und $\gamma: \Omega \rightarrow \R^n$.\\
Dann definieren wir
\[ \gamma^\sigma_k\coloneqq  \mu_{k,n}\circ \gamma \]
und nennen diese Abbildung die \textbf{$k$--te Ordnungsstatistik von $\gamma$}.
\\Ferner definieren wir
\[ \gamma_k^\ast \coloneqq  \mu_{k,n} \circ \pfeil {\abs\cdot} n \circ \gamma\]
und nennen diese Abbildung die \textbf{absolute $k$--te Ordnungsstatistik von
$\gamma$}.\\
Schließlich setzen wir:
\[ \gamma ^\sigma \coloneqq  \left(\gamma^\sigma_i \right)_{i=1}^n\quad
\text{ und }\quad \gamma^\ast \coloneqq  \left(\gamma^\ast_i \right)_{i=1}^n\]
\end{definition}